
\section{\trOrientation，\texorpdfstring{$\ZZ$}{Z} 上的\trComplex}
\label{sec:orientation-complex-over-z}

如果 $c$ 是\trMorseFunction{} $f$ 的一个\trCriticalPoint，那么它的\trStableManifold{} $W^s(c)$——它\trDiffeomorphic{}于一个\trDisk——便是一个\textbf{\trOrientable}（orientable）\trManifold。现在，我们为每一个\trCriticalPoint{} $c$ 选定\trStableManifold{} $W^s(c)$ 的一个\textbf{\trOrientation}（orientation）。请注意，这相当于选取\trUnstableManifold{} $W^u(c)$ 的一个\textbf{\trCoorientation}（co-orientation）：所谓\trVectorSpace{} $T_cW^u(c)$ 的\trCoorientation，按定义\footnote{如有需要，参见\cref{ex:ch03-14}。}就是\trComplement{} $T_cW^s(c)$ 的一个\trOrientation。

设 $a$ 与 $b$ 是两个\trCriticalPoint，于是交集 $W^u(a)\cap W^s(b)$（我们仍然假定这个交集是\trTransverse{}的）便是一个\trOrientation\trManifold。对于它与一个\trRegular\trLevelSet{}的交集，情形也是如此（\trRegular\trLevelSet{}由所采用的那个\trPseudogradient{}场 $X$ 所给出的\trTransverse\trOrientation{}赋以\trCoorientation）。于是，诸\trTrajectory{}的空间 $\mathcal{L}(a,b)$ 便是一个\trOrientation\trManifold。

我们定义 $C_{\star}(f)$ 为由 $f$ 的诸\trCriticalPoint{}所生成的自由\trAbelianGroup{}（$\ZZ$-\trModule）。让我们重复一遍：我们已经一劳永逸地在每一条 $W^s(c)$ 上固定了一个\trOrientation。于是便有
\[
C_k(f)=\Bigl\{\,\sum_{c\in\operatorname{Crit}_k(f)} a_c\,c
  \Bigm|\ a_c\in\ZZ\,\Bigr\}\text{。}
\]

当 $\operatorname{Ind}(a)=\operatorname{Ind}(b)+1$ 时，空间 $\mathcal{L}(a,b)$ 是一个维数为~0 的\trCompact\trOrientation\trManifold，也就是说，它由有限多个点组成，其中每一点都各带一个符号。我们令 $N_X(a,b)\in\ZZ$ 表示这些符号之和，并令
\[
\begin{aligned}
\partial_X\colon C_k(f)\ &\longrightarrow\ C_{k-1}(f)\\
a\ &\longmapsto\ \sum_{b\in\operatorname{Crit}_{k-1}(f)} N_X(a,b)\,b\text{。}
\end{aligned}
\]

请注意，用来定义模 2 \trComplex{}的那个 $n_X(a,b)$ 不是别的，正是 $N_X(a,b)$ 的模 2 约化。

为了保证 $\bigl(C_{\star}(f),\partial_X\bigr)$ 确实是一个\trComplex，我们还需要在这个新的语境下验证 $\partial_X\circ\partial_X=0$。\footnote{译注：原书此处印作 $\bigl(C_k(f),\partial_X\bigr)$（下标作斜体 $k$）；与本段首句的 $C_{\star}(f)$ 不一致，疑应为 $C_{\star}(f)$（显见记号笔误）；此处径改。——译者注}这当然是如下事实的一个推论：\cref{thm:broken-manifold-boundary}含有一条关于\trOrientation{}的陈述。\footnote{译注：原书此处称 \cref{thm:broken-manifold-boundary} 含有一条关于\trOrientation{}的陈述；然英文版该定理的陈述实无“\trOrientation”字样，疑法文原版之 orientée 于英译中脱落；此处照原书排印，存疑备案（所幸该陈述的结论对\trOrientation{}情形同样成立）。——译者注}显然，把一个维数为~1 的\trOrientation\trManifold\trCompactification{}成一个\trManifoldWithBoundary，所给出的便是一个\trOrientation{}的\trManifoldWithBoundary……我们尚须验证：在这个边界上诱导出的那个\trOrientation{}确实就是我们在边界上所设定的那一个。这并不是一件十分困难的事情，尤其是当我们仔细考察\cref{fig:disk-descent}或 [46,~pp.~155--156] 的时候。因此，我们把这项验证留给读者们。
